% Explicit statement inventory for the principal results in Chapters 6--7.

\section{Regular sequences and the Koszul complex}
\begin{definition}[Regular sequence]\label{mcrt-def-regular-sequence}\source{matsumura-commutative-ring-theory:page-0138}\uses{}
For an $A$-module $M$, an element $a$ is $M$-regular if $ax\ne0$ for every
$0\ne x\in M$.  A sequence $a_1,\ldots,a_r$ is an $M$-sequence when $a_1$
is $M$-regular and $a_i$ is $M/(a_1,\ldots,a_{i-1})M$-regular for each $i>1$.
\end{definition}

\begin{theorem}[Powers of a regular sequence]\label{mcrt-thm-16-1}\dcref{16.1}\source{matsumura-commutative-ring-theory:page-0138}\uses{}
If $a_1,\ldots,a_r$ is an $M$-sequence and $v_1,\ldots,v_r$ are positive
integers, then $a_1^{v_1},\ldots,a_r^{v_r}$ is an $M$-sequence.
\end{theorem}
\begin{proof}
The proof is by induction on the length and on the exponents.  If
$b_1,\ldots,b_r$ is an $M$-sequence and $\sum b_i\xi_i=0$, successively using
regularity modulo $(b_1,\ldots,b_{i-1})M$ shows
$\xi_i\in(b_1,\ldots,b_r)M$.  Applying this relation to powers gives the
required non-zero-divisor conditions for each $a_i^{v_i}$.
\end{proof}

\begin{definition}[Quasi-regular sequence]\label{mcrt-def-quasi-regular}\source{matsumura-commutative-ring-theory:page-0139}\uses{}
Put $I=(a_1,\ldots,a_r)$ and assume $IM\ne M$.  The sequence is
$M$-quasi-regular if every homogeneous $F\in M[X_1,\ldots,X_r]$ with
$F(a)\in I^{v+1}M$ has all coefficients in $IM$.
\end{definition}

\begin{theorem}[Regular implies quasi-regular]\label{mcrt-thm-16-2}\dcref{16.2}\source{matsumura-commutative-ring-theory:page-0140}\uses{mcrt-def-regular-sequence,mcrt-def-quasi-regular}
If $a_1,\ldots,a_r$ is an $M$-sequence, it is $M$-quasi-regular.  Conversely,
if $IM:x=IM$ implies $I^vM:x=I^vM$ for every $v>0$, the quasi-regular
sequence has the same coefficient criterion.
\end{theorem}
\begin{proof}
Induct on the degree of a homogeneous polynomial.  For the converse, write a
coefficient relation in degree $v$ as $xF(a)\in I^{v+1}M$ and apply the
induction hypothesis followed by $IM:x=IM$.  For the forward implication,
separate the last variable and use regularity modulo the preceding terms;
the resulting coefficient relation reduces to the induction hypothesis.
\end{proof}

\begin{theorem}[Quasi-regular implies regular]\label{mcrt-thm-16-3}\dcref{16.3}\source{matsumura-commutative-ring-theory:page-0141}\uses{mcrt-thm-16-2}
Let $A$ be Noetherian, $M\ne0$, and suppose that $M,M/a_1M,\ldots,
M/(a_1,\ldots,a_{r-1})M$ are $I$-adically separated.  Then every
$M$-quasi-regular sequence $a_1,\ldots,a_r$ is an $M$-sequence.
\end{theorem}
\begin{proof}
If $a_1x=0$, separatedness and the quasi-regular coefficient criterion force
$x\in I^vM$ for all $v$, hence $x=0$.  Apply the same argument to
$M/a_1M$ and induct on $r$.
\end{proof}

\begin{definition}[Koszul complex]\label{mcrt-def-koszul}\source{matsumura-commutative-ring-theory:page-0142}\uses{}
For $x_1,\ldots,x_r\in A$, let $K_p$ be free on symbols
$e_{i_1}\wedge\cdots\wedge e_{i_p}$ and define
$d(e_{i_1}\wedge\cdots\wedge e_{i_p})=\sum_j(-1)^{j-1}x_{i_j}
e_{i_1}\wedge\widehat e_{i_j}\wedge\cdots\wedge e_{i_p}$.  Write
$K_\bullet(x,M)=K_\bullet(x)\otimes_A M$ and $H_p(x,M)$ for its homology.
\end{definition}

\begin{theorem}[Koszul homology criterion]\label{mcrt-thm-16-5}\dcref{16.5}\source{matsumura-commutative-ring-theory:page-0143}\uses{mcrt-def-koszul,mcrt-def-regular-sequence}
If $x_1,\ldots,x_r$ is an $M$-sequence, then $H_p(x,M)=0$ for $p>0$ and
$H_0(x,M)=M/(x_1,\ldots,x_r)M$.  Under the local finite or positively graded
hypotheses, the converse holds when $M\ne0$.
\end{theorem}
\begin{proof}
Induct on $r$ using the short exact sequence obtained by tensoring with the
two-term Koszul complex on the last element.  The long exact homology sequence
and regularity give vanishing in positive degrees.  Conversely, Nakayama's
lemma applied to the same sequence gives regularity one element at a time.
\end{proof}

\begin{theorem}[Ext characterization of depth]\label{mcrt-thm-16-6}\dcref{16.6}\source{matsumura-commutative-ring-theory:page-0144}\uses{mcrt-thm-16-3}
For a Noetherian ring $A$, finite $M$ with $IM\ne M$, and $n>0$, the following
are equivalent: $\operatorname{Ext}^i_A(N,M)=0$ for all finite $N$ and $i<n$;
the same vanishing for $N=A/I$; and the existence of an $M$-sequence of length
$n$ in $I$.
\end{theorem}
\begin{proof}
The implication from a regular sequence follows by induction on its length and
the long exact Ext sequence.  For the converse, if every element of $I$ is a
zero-divisor, an associated prime containing $I$ produces a nonzero
$\operatorname{Hom}(A/I,M)$, contradicting the vanishing.  Choose a regular
element and repeat on the quotient.
\end{proof}

\begin{theorem}[Depth]\label{mcrt-thm-16-7}\dcref{16.7}\source{matsumura-commutative-ring-theory:page-0145}\uses{mcrt-thm-16-6}
The length of a maximal $M$-sequence in $I$ is a well-determined integer
$\operatorname{depth}(I,M)$, characterized by
$\operatorname{Ext}^i_A(A/I,M)=0$ for $i<\operatorname{depth}(I,M)$ and
nonvanishing in the depth degree.
\end{theorem}
\begin{proof}
Extend any $M$-sequence while the next Ext group vanishes, using the preceding
theorem.  A maximal sequence therefore stops exactly at the first nonzero Ext
group, so all maximal sequences have the same length.
\end{proof}

\begin{theorem}[Koszul depth sensitivity]\label{mcrt-thm-16-8}\dcref{16.8}\source{matsumura-commutative-ring-theory:page-0146}\uses{mcrt-thm-16-5,mcrt-thm-16-7}
If $I=(y_1,\ldots,y_n)$ and $M\ne IM$, then a maximal $M$-sequence in $I$
has length $n-q$, where $q=\sup\{i:H_i(y,M)\ne0\}$.
\end{theorem}
\begin{proof}
Induct on the length of a maximal sequence and use the long exact sequence for
the first regular element.  The Koszul homology of the quotient shifts by one,
which gives $q=n-s$ when the sequence has length $s$.
\end{proof}

\begin{theorem}[Grade versus projective dimension]\label{mcrt-thm-16-9}\dcref{16.9}\source{matsumura-commutative-ring-theory:page-0147}\uses{mcrt-thm-16-7}
If $M,N$ are finite Noetherian modules, $M\ne0$, $\operatorname{grade}M=k$,
and $\operatorname{projdim}N=\ell\le k$, then
$\operatorname{Ext}^i_A(M,N)=0$ for $i<k-\ell$.
\end{theorem}
\begin{proof}
Induct on $\ell$.  Resolve $N$ by a finite free module and use the long exact
Ext sequence; the case $\ell=0$ is the definition of grade.
\end{proof}

\section{Cohen--Macaulay rings}
\begin{definition}[Cohen--Macaulay module]\label{mcrt-def-cm}\source{matsumura-commutative-ring-theory:page-0149}\uses{}
A finite module over a Noetherian local ring is Cohen--Macaulay if it is zero or
if $\depth M=\dim M$.  The ring is Cohen--Macaulay when it is so as a module.
\end{definition}

\begin{theorem}[Ischebeck]\label{mcrt-thm-17-1}\dcref{17.1}\source{matsumura-commutative-ring-theory:page-0148}\uses{mcrt-thm-16-6}
For nonzero finite modules over a Noetherian local ring, if $\depth M=k$ and
$\dim N=r$, then $\operatorname{Ext}^i_A(N,M)=0$ for $i<k-r$.
\end{theorem}
\begin{proof}
Induct on $r$, filter $N$ by cyclic modules $A/P$, and for $r>0$ choose
$x\notin P$ to reduce to $A/(P,x)$; the long exact sequence and Nakayama's
lemma finish the induction.
\end{proof}

\begin{theorem}[Associated-prime dimension bound]\label{mcrt-thm-17-2}\dcref{17.2}\source{matsumura-commutative-ring-theory:page-0149}\uses{mcrt-thm-17-1}
If $P\in\operatorname{Ass}(M)$, then $\dim(A/P)\ge\depth M$.
\end{theorem}
\begin{proof}
There is a nonzero map $A/P\to M$.  Apply the preceding Ext vanishing theorem;
otherwise $\operatorname{Hom}(A/P,M)$ would be forced to vanish.
\end{proof}

\begin{theorem}[Cohen--Macaulay module properties]\label{mcrt-thm-17-3}\dcref{17.3}\source{matsumura-commutative-ring-theory:page-0149}\uses{mcrt-def-cm,mcrt-thm-17-2,mcrt-thm-16-7}
For a finite module $M$ over a Noetherian local ring: (i) a CM module has no
embedded associated primes; (ii) quotienting by an $M$-sequence preserves CM;
(iii) localization of a CM module is CM.
\end{theorem}
\begin{proof}
Use the dimension inequalities from the associated-prime bound.  Quotienting
by a regular sequence lowers both depth and dimension by its length; localizing
and inducting on localized depth gives the third assertion.
\end{proof}

\begin{theorem}[Cohen--Macaulay criterion]\label{mcrt-thm-17-4}\dcref{17.4}\source{matsumura-commutative-ring-theory:page-0150}\uses{mcrt-thm-17-3,mcrt-thm-16-7}
For a CM local ring $A$ and a proper ideal $I$,
$\operatorname{ht}I=\depth(I,A)=\operatorname{grade}I$ and
$\operatorname{ht}I+\dim A/I=\dim A$.  For $a_1,\ldots,a_r\in\mathfrak m$,
the conditions that they form an $A$-sequence, have successive height $i$,
generate height $r$, or form part of a parameter system are equivalent.
\end{theorem}
\begin{proof}
Apply the preceding theorem to parameter systems and localizations.  Induct on
$\dim A$ to obtain regularity of every parameter system, then compare heights
of minimal prime divisors to get the dimension formula and the equivalences.
\end{proof}

\begin{theorem}[Completion and Cohen--Macaulayness]\label{mcrt-thm-17-5}\dcref{17.5}\source{matsumura-commutative-ring-theory:page-0151}\uses{mcrt-thm-16-7}
For a Noetherian local ring $A$ and its completion $\widehat A$,
$\depth A=\depth\widehat A$ and $A$ is CM if and only if $\widehat A$ is CM.
\end{theorem}
\begin{proof}
Compare $\operatorname{Ext}^i_A(k,A)$ with the corresponding Ext groups after
completion and use invariance of dimension under completion.
\end{proof}

\begin{theorem}[Unmixedness criterion]\label{mcrt-thm-17-6}\dcref{17.6}\source{matsumura-commutative-ring-theory:page-0151}\uses{mcrt-thm-17-4}
A Noetherian ring is CM if and only if every ideal of height $r$ generated by
$r$ elements is unmixed.
\end{theorem}
\begin{proof}
Localize an embedded prime and use the CM criterion to contradict embeddedness.
Conversely, choose elements successively with the required heights; unmixedness
then makes each next element regular after localization.
\end{proof}

\begin{theorem}[Polynomial extension]\label{mcrt-thm-17-7}\dcref{17.7}\source{matsumura-commutative-ring-theory:page-0152}\uses{mcrt-thm-17-4}
If $A$ is a CM ring, then $A[X_1,\ldots,X_n]$ is CM.
\end{theorem}
\begin{proof}
Reduce to one variable and localize at a maximal ideal.  Lift the monic
irreducible residue polynomial and append it to a parameter sequence; flatness
preserves the regular sequence and gives the dimension bound.
\end{proof}

\begin{theorem}[Regular local rings are CM]\label{mcrt-thm-17-8}\dcref{17.8}\source{matsumura-commutative-ring-theory:page-0152}\uses{mcrt-thm-17-4}
Every regular local ring is Cohen--Macaulay.
\end{theorem}
\begin{proof}
A regular system of parameters generates a strictly increasing chain of prime
ideals, hence is a regular sequence by the height criterion.
\end{proof}

\begin{theorem}[Quotients of CM rings are universally catenary]\label{mcrt-thm-17-9}\dcref{17.9}\source{matsumura-commutative-ring-theory:page-0152}\uses{mcrt-thm-17-4}
Every quotient of a Cohen--Macaulay ring is universally catenary.
\end{theorem}
\begin{proof}
This follows by applying the dimension formula in the CM criterion to all
localizations and quotients.
\end{proof}

\begin{theorem}[Graded characterization of regularity]\label{mcrt-thm-17-10}\dcref{17.10}\source{matsumura-commutative-ring-theory:page-0153}\uses{mcrt-thm-16-5,mcrt-thm-17-4}
A Noetherian local ring $(A,\mathfrak m,k)$ is regular if and only if
$\operatorname{gr}_{\mathfrak m}(A)$ is a polynomial $k$-algebra.
\end{theorem}
\begin{proof}
For a regular system of parameters, the Koszul/graded calculation identifies
the associated graded ring with a polynomial ring.  Conversely, compare the
degree-one and degree-$n$ dimensions to obtain
$\dim A=\dim_k\mathfrak m/\mathfrak m^2$.
\end{proof}

\begin{theorem}[Multiplicity criterion]\label{mcrt-thm-17-11}\dcref{17.11}\source{matsumura-commutative-ring-theory:page-0153}\uses{mcrt-thm-17-10,mcrt-thm-17-4}
For a Noetherian local ring, the following are equivalent: $A$ is CM;
$\length(A/q)=e(q)$ for every parameter ideal $q$; and this equality holds for
some parameter ideal.
\end{theorem}
\begin{proof}
For a parameter sequence compute the Hilbert function from the graded free
polynomial extension.  Equality of multiplicity and length forces the graded
kernel to vanish, which is equivalent to the CM condition.
\end{proof}

\section{Gorenstein rings}
\begin{definition}[Gorenstein local ring]\label{mcrt-def-gorenstein}\source{matsumura-commutative-ring-theory:page-0157}\uses{}
A Noetherian local ring is Gorenstein when its injective dimension as a module
over itself is finite (equivalently, its top Bass number is one).
\end{definition}
\begin{theorem}[Gorenstein permanence]\label{mcrt-thm-18-1}\dcref{18.1}\source{matsumura-commutative-ring-theory:page-0153}\uses{mcrt-def-gorenstein,mcrt-thm-17-8}
For a Noetherian local ring, the standard equivalent Gorenstein conditions
include finite self-injective dimension, one-dimensional socle in the top Ext,
and the corresponding Bass-number and canonical-module criteria.  Regular local
rings are Gorenstein, and quotients by regular elements preserve the property.
\end{theorem}
\begin{proof}
Use a minimal injective resolution and the Ext characterization of depth.  The
long exact sequence for a regular element identifies the shifted Bass numbers,
and induction gives the quotient and regular-local assertions.
\end{proof}

\begin{theorem}[Gorenstein localization]\label{mcrt-thm-18-2}\dcref{18.2}\source{matsumura-commutative-ring-theory:page-0159}\uses{mcrt-thm-18-1}\notready
If $A$ is a Gorenstein local ring and $P\in\operatorname{Spec}A$, then $A_P$
is Gorenstein.
\end{theorem}
\begin{proof}
Localize the finite injective resolution and compare the resulting Bass numbers
at $P$.
\end{proof}
